\section{Topología dinámica: qué se rompe y cómo se repara}
\label{sec:topologia}

Los resultados de consenso sobre multiplicadores y sobre \emph{log-odds}
invocados hasta aquí suponen un grafo fijo: fuertemente conexo y balanceado en
peso en el caso de los multiplicadores \cite{martinezPiazuelo2025gne}, no dirigido
y conexo en \cref{th:mapa}. Esa hipótesis es insostenible en este problema,
porque el grafo de comunicación es el de disco
$\Gcal_{\mathrm{disk}}(R^{\mathrm{com}})$ y los AMR se mueven: acercarse o
alejarse crea y destruye aristas.

Tres resultados dependen de ello y conviene decir exactamente cuáles.
\Cref{th:mapa} usa $\lambda_2$ de un Laplaciano fijo. \Cref{th:precio-seguridad}
exige que los dos miembros de un par \emph{acuerden} el precio, y el acuerdo se
obtiene por consenso sobre el mismo grafo. Y la noción de desviación
\emph{conectada} de \cref{sec:nhop}, que define $h_c^\star$, se enuncia sobre un
grafo que ahora varía. Se presentan dos salidas, y la segunda es la que se
adopta.

\subsection{Salida 1: tolerar la conmutación}

La teoría clásica de consenso bajo topología conmutada no exige conectividad
instantánea, sino acumulada.

\begin{teorema}[Consenso bajo topología conmutada]
\label{th:conmutada}
Sea la topología constante a trozos en instantes $t_0<t_1<\cdots$ con tiempos
de permanencia $\tau_j=t_{j+1}-t_j$ acotados inferiormente por $\tau_{\min}>0$.
Si la unión de los grafos dirigidos sobre intervalos contiguos acotados posee un
árbol generador dirigido, entonces el protocolo de consenso alcanza acuerdo
\cite[Lemas 2.26--2.27]{ren2008distributed}.
\end{teorema}

La demostración clásica pasa por productos infinitos de matrices SIA
(estocásticas, indescomponibles y aperiódicas): cada $e^{-L(t_j)\tau_j}$ es
estocástica por filas, y la unión con árbol generador garantiza que el producto
lo sea también.

\begin{observacion}[Lo que cuesta esta salida]
\label{obs:coste-conmutada}
La tasa deja de ser $(1-\varepsilon\lambda_2)^t$ con un $\lambda_2$ fijo: pasa a
depender de la longitud del intervalo sobre el que hay que acumular el árbol
generador, y se degrada cuando la conectividad es intermitente. Además, la
hipótesis es sobre la \emph{trayectoria realizada}: no puede verificarse
\emph{a priori} porque el grafo lo produce el propio movimiento que el
mecanismo decide. Es una hipótesis sobre el resultado, no sobre el diseño.
\end{observacion}

\subsection{Salida 2: garantizar la topología con una barrera}

La circularidad de \cref{obs:coste-conmutada} se rompe si se \emph{impone} la
conectividad en lugar de suponerla. La observación es que una arista de
comunicación es una restricción de distancia, y por tanto admite exactamente el
mismo tratamiento que la seguridad en \cref{sec:clusters}: una barrera con su
multiplicador.

Sea $\mathcal T\subseteq\mathcal E_C(t_0)$ un árbol generador del grafo de
comunicación en el instante inicial y, para cada $(i,j)\in\mathcal T$,
\begin{equation}
  h^{\mathrm{con}}_{ij}
  \defeq (R^{\mathrm{com}})^2-\|p_i-p_j\|^2\ \ge\ 0 .
  \label{eq:barrera-conectividad}
\end{equation}

\begin{teorema}[Topología garantizada]
\label{th:topologia-garantizada}
Si $h^{\mathrm{con}}_{ij}(t_0)\ge0$ para todo $(i,j)\in\mathcal T$ y el QP de
movimiento que incluye las barreras \eqref{eq:barrera-conectividad} es factible
en todo estado alcanzado, entonces $\mathcal T\subseteq\mathcal E_C(t)$ para
todo $t\ge t_0$. En consecuencia el grafo de comunicación permanece conexo, y
las hipótesis de grafo fijo de \cref{th:mapa} y de la dinámica de
multiplicadores se satisfacen sobre $\mathcal T$.
\end{teorema}

\begin{proof}
Cada $h^{\mathrm{con}}_{ij}$ es una función de barrera con grado relativo dos
respecto a la aceleración, idéntica en estructura a la de seguridad. Aplicando
\cref{pr:hocbf} al conjunto $\{h^{\mathrm{con}}_{ij}\ge0\}_{(i,j)\in\mathcal T}$
se obtiene invariancia hacia adelante, luego ninguna arista de $\mathcal T$ se
rompe. Un grafo que contiene un árbol generador es conexo.
\end{proof}

\Cref{th:topologia-garantizada} convierte un problema de topología conmutada en
uno de topología fija \emph{por construcción}. Las aristas que no pertenecen a
$\mathcal T$ aparecen y desaparecen libremente; solo añaden información y su
pérdida no compromete ningún resultado.

\subsection{Seguridad y conectividad están en tensión}

Las dos barreras empujan en sentidos opuestos: la de seguridad separa, la de
conectividad acerca. Su compatibilidad no es gratuita.

\begin{proposicion}[Condición de compatibilidad y cota de grado]
\label{pr:compatibilidad}
Para un par, el conjunto admisible es el anillo
$d_{\min}\le\|p_i-p_j\|\le R^{\mathrm{com}}$, no vacío si y solo si
$d_{\min}\le R^{\mathrm{com}}$. Para un nodo de grado $\Delta$ en $\mathcal T$,
una condición necesaria es
\[
  \Delta\ \le\
  \left\lfloor
    \frac{\pi}{\arcsin\!\bigl(d_{\min}/(2R^{\mathrm{com}})\bigr)}
  \right\rfloor .
\]
\end{proposicion}

\begin{proof}
Los $\Delta$ vecinos deben estar a distancia no mayor que $R^{\mathrm{com}}$ del
nodo y separados entre sí al menos $d_{\min}$. La disposición más permisiva los
sitúa sobre la circunferencia de radio $R^{\mathrm{com}}$: dos vecinos a
separación angular $\theta$ distan $2R^{\mathrm{com}}\sin(\theta/2)$, que debe
ser al menos $d_{\min}$, luego
$\theta\ge2\arcsin\bigl(d_{\min}/(2R^{\mathrm{com}})\bigr)$. Caben a lo sumo
$\lfloor 2\pi/\theta\rfloor$ vecinos. Situarlos a radio menor exige separación
angular mayor, luego la cota es necesaria.
\end{proof}

La cota es operativa: fija qué árboles generadores son admisibles. Con
$d_{\min}=0{,}5$~m y $R^{\mathrm{com}}=3{,}2$~m resulta $\Delta\le40$, holgado;
con $d_{\min}=1{,}0$~m y $R^{\mathrm{com}}=1{,}5$~m resulta $\Delta\le9$, que ya
condiciona la elección del árbol.

\begin{corolario}[El precio de la conectividad]
\label{co:precio-conectividad}
El multiplicador de $h^{\mathrm{con}}_{ij}$ es el coste marginal de mantener el
enlace $(i,j)$. Comparado con el valor de la información que transporta,
proporciona un criterio para reconfigurar $\mathcal T$: cuando el precio de una
arista supera de forma sostenida ese valor, conviene recablear el árbol. La
reconfiguración es una conmutación estructural y queda sujeta a
\cref{th:dwell}.
\end{corolario}

\begin{observacion}[Qué se paga y dónde aparece]
\label{obs:coste-conectividad}
Imponer $\mathcal T$ restringe el movimiento: hay rutas admisibles sin la
restricción que dejan de serlo con ella. Esa pérdida no desaparece, se
contabiliza: entra en $\Delta_{\mathrm{dyn}}$ de \cref{th:descomposicion}. La
elección entre las dos salidas es, por tanto, explícita. La primera no restringe
el movimiento y ofrece una garantía condicionada a una hipótesis que solo puede
comprobarse a posteriori; la segunda restringe el movimiento y ofrece una
garantía verificable en diseño.
\end{observacion}

\begin{observacion}[Atribución]
\label{obs:atribucion-conectividad}
El mantenimiento de conectividad mediante campos potenciales o barreras es
técnica conocida en control multi-robot; lo que aquí se sostiene no es la
técnica sino su lugar en la arquitectura: la conectividad es una barrera más y,
por \cref{th:precio-seguridad}, su multiplicador es un precio más dentro del
mismo juego. No se introduce ningún mecanismo adicional.
\end{observacion}
